\label{sec:production}
% MAIN TEXT, v8 Section 2 (Model).
% In the merged Model section, replace the material after eq:budget as
% specified in PLACEMENT_KO.md. Keep the earlier workload accounting.

The planning horizon and the resource numeraire are fixed. Let
$B_{\rm res}$ be physical resources available during that horizon and let
$b_{\rm ref}>0$ be a fixed reference amount. The budget coordinate is
$B=B_{\rm res}/b_{\rm ref}$; local investment and project prices use the
same resource unit. Thus $p_jx_j$ and $z$ are resource expenditures, whereas
$x,y,w$ are dimensionless increments of project capacity. Similarly,
$C=C_{\rm useful}/c_{\rm ref}$ uses a fixed computational-service unit
$c_{\rm ref}>0$. In particular, logarithms take dimensionless arguments.
An efficiency improvement changes the production coefficient $a$, not
$B$, unless it also changes the physical resources actually available.
Other physical services retain the separate constraints in \cref{eq:chi}.

The cumulative indices $n=1+x$, $h=1+x+y$, and $m=1+x+y+w$ retain the
residence, industry, and observation accounting of the model. Useful
computation and utility are
\begin{align}
 C&=C_0+a\ell(z)+G_n(x)+G_h(y),\label{eq:compute}\\
 U&=F(C)+(1-q)E(x)+\epsilon_RR(x+y)
       +\epsilon_II(x+y+w).\label{eq:utility}
\end{align}
Here $a>0$, $C_0\geq0$, $q\in[0,1]$, and the residual weights are
non-negative. The functions $E,R,I$ are differentiable, non-decreasing,
and concave on the feasible domain. The functions $G_n,G_h$ are
continuous, non-decreasing, and concave, with $G_j(0)=0$ and finite right
derivatives $G_j'(0)=r_j\geq0$. Set $r_m=0$. Local production is continuous
and non-negative on $[0,B]$, with $\ell(0)=0$, and has a positive finite
left derivative at $B$ and $\ell(B)>0$. Whenever global concavity is
invoked, $\ell$ is also concave on this interval. These assumptions do
not require a derivative at zero. The baseline is
$\ell(z)=\ln(1+z)$ and $G_j(v)=r_jv$.

We use the shifted-isoelastic computational value
\begin{equation}
 F_\eta(C)=
 \begin{cases}
 [(1+C)^{1-\eta}-1]/(1-\eta),&\eta\ne1,\\
 \ln(1+C),&\eta=1,
 \end{cases}
 \qquad \eta\geq0,
 \label{eq:isoelastic}
\end{equation}
so that $F_\eta'(C)=(1+C)^{-\eta}$, $F_\eta(0)=0$, and
$F_\eta'(0)=1$. The common marginal-value normalization makes the
residual weights comparable across this family. The value-curvature
parameter $\eta$ is distinct from the service-efficiency factors $\eta_k$.
Recursive description alone does not determine a welfare functional.
% Insert the Stage 3 free-world-number lemma here after it is established.

The residual slopes at the all-inward corner are
\begin{align}
 S_n&=(1-q)E'(0)+\epsilon_RR'(0)+\epsilon_II'(0),\nonumber\\
 S_h&=\epsilon_RR'(0)+\epsilon_II'(0),&
 S_m&=\epsilon_II'(0).
 \label{eq:slopes}
\end{align}
Write $d_0=1+C_0\geq1$ and define
\begin{equation}
 \Omega=a\ell(B),\qquad
 \Gamma(B)=\frac{\ell(B)}{\ell'(B)}.
 \label{eq:omega-gamma}
\end{equation}
The first quantity is new local computational output at the corner;
the second is a resource-scale measure of production curvature.
Changing the multiplicative normalization of $\ell$ and compensating in
$a$ changes neither quantity. Throughout, all-inward allocation means
zero new expenditure in every external direction. It is stronger than
zero new replication, since observation can be external without
replicating.
